\section{Notação e Propriedades do MCS}

Esta seção apresenta a notação básica e as propriedades do MCS.

\begin{itemize}

    %%%%%%%%%%%%%%%%%%%%%%%%%%%%%
    % Basic objects and notation
    %%%%%%%%%%%%%%%%%%%%%%%%%%%%%

  \item \textbf{Basic sets and losses}:
    \begin{itemize}
      \item Initial set of models (objects): $M_0 = \{1,\dots,m_0\}$. [file:1]
      \item Loss of model $i$ at time $t$: $L_{i,t}$ for $t=1,\dots,n$. [file:1]
      \item Loss differential between models $i$ and $j$:
        \[
          d_{ij,t} \equiv L_{i,t} - L_{j,t}, \quad i,j \in M_0.
        \] [file:1]
      \item Expected loss differential:
        \[
          \mu_{ij} \equiv \mathbb{E}(d_{ij,t}), \quad \mu_{ij} \text{ does not depend on } t.
        \] [file:1]
      \item Model $i$ is preferred to $j$ if $\mu_{ij} < 0$. [file:1]
    \end{itemize}

  \item \textbf{Set of superior models (Definition 1)}:
    \begin{itemize}
      \item The set of superior objects is
        \[
          M^\ast \equiv \bigl\{ i \in M_0 : \mu_{ij} \le 0 \ \text{for all } j \in M_0 \bigr\}.
        \] [file:1]
      \item $M^\ast$ may contain one or multiple models. [file:1]
    \end{itemize}

  \item \textbf{Null and alternative hypotheses}:
    \begin{itemize}
      \item For any subset $M \subseteq M_0$, define the null of equal performance
        \[
          H_{0,M} : \mu_{ij} = 0 \quad \text{for all } i,j \in M.
        \] [file:1]
      \item The corresponding alternative is
        \[
          H_{A,M}: \mu_{ij} \neq 0 \ \text{for some } i,j \in M.
        \] [file:1]
      \item By definition, $H_{0,M^\ast}$ is true; $H_{0,M}$ is false if $M$ contains elements both from $M^\ast$ and from $M_0 \setminus M^\ast$. [file:1]
    \end{itemize}

    %%%%%%%%%%%%%%%%%%%%%%%%%%%%%
    % MCS algorithm and objects
    %%%%%%%%%%%%%%%%%%%%%%%%%%%%%

  \item \textbf{Equivalence test and elimination rule}:
    \begin{itemize}
      \item For each $M \subseteq M_0$, an equivalence test $\delta_M \in \{0,1\}$ is used to test $H_{0,M}$:
        \[
          \delta_M = 0 \iff \text{``accept'' } H_{0,M}, \quad
          \delta_M = 1 \iff \text{reject } H_{0,M}.
        \] [file:1]
      \item For each $M \subseteq M_0$, an elimination rule $e_M \in M$ selects one model to be removed from $M$ when $H_{0,M}$ is rejected. [file:1]
    \end{itemize}

  \item \textbf{MCS algorithm (Definition 2)}:
    \begin{itemize}
      \item \emph{Step 0}: Initialize $M = M_0$. [file:1]
      \item \emph{Step 1}: Test $H_{0,M}$ using $\delta_M$ at significance level $\alpha$. [file:1]
      \item \emph{Step 2}:
        \begin{itemize}
          \item If $H_{0,M}$ is accepted ($\delta_M=0$), define the model confidence set
            \[
              \widehat{M}^\ast_{1-\alpha} = M.
            \]
            and stop. [file:1]
          \item If $H_{0,M}$ is rejected ($\delta_M=1$), remove $e_M$ from $M$ and return to Step 1. [file:1]
        \end{itemize}
      \item The resulting random set $\widehat{M}^\ast_{1-\alpha}$ of surviving models is called the \emph{model confidence set (MCS)} at confidence level $1-\alpha$. [file:1]
    \end{itemize}

  \item \textbf{Asymptotic assumptions on $(\delta_M,e_M)$ (Assumption 1)}:
    \begin{itemize}
      \item For each $M \subseteq M_0$:
        \begin{align*}
          &\text{(a) } \limsup_{n\to\infty} \mathbb{P}(\delta_M = 1 \mid H_{0,M}) \le \alpha &&\text{(asymptotic level $\le\alpha$),} \\
          &\text{(b) } \lim_{n\to\infty} \mathbb{P}(\delta_M = 1 \mid H_{A,M}) = 1 &&\text{(asymptotic power $=1$),} \\
          &\text{(c) } \lim_{n\to\infty} \mathbb{P}(e_M \in M^\ast \mid H_{A,M}) = 0 &&\text{(do not eliminate superior models asymptotically).}
        \end{align*} [file:1]
    \end{itemize}

    %%%%%%%%%%%%%%%%%%%%%%%%%%%%%
    % Theorem 1
    %%%%%%%%%%%%%%%%%%%%%%%%%%%%%

  \item \textbf{Theorem 1 (Properties of the MCS)}:
    \begin{itemize}
      \item Under Assumption 1, the MCS has the following asymptotic coverage and exclusion properties:
        \begin{align*}
          &\text{(i)} \quad \liminf_{n\to\infty} \mathbb{P}\bigl( M^\ast \subseteq \widehat{M}^\ast_{1-\alpha} \bigr) \;\ge\; 1 - \alpha, \\[0.4em]
          &\text{(ii)} \quad \lim_{n\to\infty} \mathbb{P}\bigl( i \in \widehat{M}^\ast_{1-\alpha} \bigr) = 0 \quad \text{for all } i \notin M^\ast.
        \end{align*} [file:1]
      \item (i) says the MCS covers all superior models with asymptotic probability at least $1-\alpha$. [file:1]
      \item (ii) says any inferior model is asymptotically excluded from the MCS. [file:1]
    \end{itemize}

    %%%%%%%%%%%%%%%%%%%%%%%%%%%%%
    % Corollary 1
    %%%%%%%%%%%%%%%%%%%%%%%%%%%%%

  \item \textbf{Corollary 1 (Singleton $M^\ast$)}:
    \begin{itemize}
      \item If $M^\ast$ is a singleton, say $M^\ast = \{i^\ast\}$, and Assumption 1 holds, then
        \[
          \lim_{n\to\infty} \mathbb{P}\bigl( M^\ast = \widehat{M}^\ast_{1-\alpha} \bigr) = 1.
        \] [file:1]
      \item Thus, when there is a unique best model, the MCS procedure selects exactly that model with asymptotic probability 1. [file:1]
    \end{itemize}

    %%%%%%%%%%%%%%%%%%%%%%%%%%%%%
    % Coherency and Theorem 2
    %%%%%%%%%%%%%%%%%%%%%%%%%%%%%

  \item \textbf{Coherency between test and elimination rule (Definition 3)}:
    \begin{itemize}
      \item Let $P$ denote the true probability measure and $P_0$ the measure obtained by imposing the null via the transformation $d_{ij,t} \mapsto d_{ij,t} - \mu_{ij}$, so that $P_0$ satisfies $H_{0,M}$. [file:1]
      \item The test $\delta_M$ and elimination rule $e_M$ are said to be \emph{coherent} if
        \[
          \mathbb{P}\bigl( \delta_M = 1,\ e_M \in M^\ast \bigr) \;\le\; P_0(\delta_M = 1).
        \] [file:1]
      \item Intuitively, coherence controls the probability of eliminating a superior model by relating it to the null rejection probability. [file:1]
    \end{itemize}

  \item \textbf{Theorem 2 (Finite-sample coverage under coherency)}:
    \begin{itemize}
      \item Suppose
        \[
          \mathbb{P}\bigl( \delta_M = 1,\ e_M \in M^\ast \bigr) \le \alpha
        \]
        for every $M \subseteq M_0$. [file:1]
      \item Then the MCS has finite-sample coverage
        \[
          \mathbb{P}\bigl( M^\ast \subseteq \widehat{M}^\ast_{1-\alpha} \bigr) \ge 1 - \alpha.
        \] [file:1]
      \item Thus, if coherency bounds the probability of wrongly eliminating a superior model by $\alpha$, the MCS is a finite-sample $(1-\alpha)$ confidence set for $M^\ast$. [file:1]
    \end{itemize}

    %%%%%%%%%%%%%%%%%%%%%%%%%%%%%
    % MCS p-values and Theorem 3
    %%%%%%%%%%%%%%%%%%%%%%%%%%%%%

  \item \textbf{Elimination sequence and MCS p-values (Definition 4)}:
    \begin{itemize}
      \item The elimination rule $e_M$ induces an ordered sequence of sets
        \[
          M_0 = M_1 \supset M_2 \supset \cdots \supset M_{m_0},
        \]
        where each $M_j$ is obtained from $M_{j-1}$ by removing $e_{M_{j-1}}$, and
        \[
          M_j = \{ e_{M_j}, e_{M_{j+1}}, \dots, e_{M_{m_0}} \}.
        \] [file:1]
      \item Let $P_{H_{0,M_i}}$ denote the p-value associated with testing $H_{0,M_i}$; set by convention $P_{H_{0,M_{m_0}}} \equiv 1$. [file:1]
      \item For a model $e_{M_j} \in M_0$, define its MCS p-value as
        \[
          \widehat{p}_{e_{M_j}} \equiv \max_{i \le j} P_{H_{0,M_i}}.
        \] [file:1]
      \item This construction yields a nondecreasing sequence
        \[
          \widehat{p}_{e_{M_1}} \le \widehat{p}_{e_{M_2}} \le \cdots \le \widehat{p}_{e_{M_{m_0}}}.
        \] [file:1]
    \end{itemize}

  \item \textbf{Theorem 3 (Characterization of MCS via p-values)}:
    \begin{itemize}
      \item Index the elements of $M_0$ by $i=1,\dots,m_0$. For each model $i$, let $\widehat{p}_i$ be its MCS p-value as defined above. [file:1]
      \item Then, for any significance level $\alpha$,
        \[
          i \in \widehat{M}^\ast_{1-\alpha}
          \quad \Longleftrightarrow \quad
          \widehat{p}_i \ge \alpha.
        \] [file:1]
      \item Equivalently, a model belongs to the $(1-\alpha)$ MCS if and only if its MCS p-value is at least $\alpha$. [file:1]
    \end{itemize}

\end{itemize}